% ============================================================
% Chapter 5: Development
%   Section: News Acquisition: Discovery and Extraction
%     Subsection: Verification Against the Live Sources
% Included from chapters/05-development/02-news-acquisition/news-acquisition.tex
% ============================================================

\subsection{Verification Against the Live Sources}
\label{subsec:scraping-live}

Unit tests show that the code does what it was written to do; only the
real sources show whether that is enough. The acquisition layer was
checked against them at each stage of its construction
(Table~\ref{tab:scraping-live}).

\begin{table}[ht]
\centering
\small
\caption{Checks of the acquisition layer against the live sources
(2026).}
\label{tab:scraping-live}
\begin{tabularx}{\linewidth}{@{}lp{1.5cm}X@{}}
\toprule
Date & Sources & Result \\
\midrule
23 Sep & 12
& Discovery went from 4 sources producing links to 9. One NASA article
  went from discovery to stored analysis in 52~s, and the next run
  reported it as already stored. \\
24 Sep & 12
& Source check: 6 working, every sample with title, author and date;
  6 broken (five discover nothing; El País gives 147 links and answers
  403 to every article, headless browser included). \\
25 Sep & 12
& First probe: 6 up, 3 degraded (dead feeds rescued by section pages:
  113, 115 and 64 links, 6 of 6 articles each), 3 down (EFE 403,
  Reuters 401, El País 403 on articles). Feeds 395 links, section pages
  1,272 more; 53 of 60 articles extracted. \\
29 Sep & 34
& Probe: 29 up, 4 degraded, 1 down; 156 of 165 sampled articles
  extracted (95\%). Feeds 1,314 links, section pages 3,604 more. \\
4 Oct  & 29 feeds
& All 1,248 discovered URLs carry a publication time. \\
6 Oct  & 2
& EFE and Reuters disabled: in the container's statistics since
  23~September, 21 and 24 attempts, none delivered. \\
\bottomrule
\end{tabularx}
\end{table}

Three observations follow from these runs.

\paragraph{Section pages are not a fallback of last resort.} On both
probes the section pages found several times more links than the feeds
(1,272 against 395, then 3,604 against 1,314). They are still tried
last, because a feed gives a time, a title and a summary for every
link, and a section page gives none of them.

\paragraph{What fails is access, not parsing.} Every article that
extracted on 25~September was extracted by trafilatura, the first
step. The failures were refusals (403, 401) and dead feeds, which is
why the cascade spends its effort on deciding when to stop rather than
on more parsers. Of the 29~September failures, El País gave 2 of 5
(the other three answered 403, and that host had no Chromium to
escalate to), Reuters refused everything, and Positive News' two
failures were section pages taken for articles, not failures of the
site.

\paragraph{Diagnoses can be wrong in instructive ways.} The four dead
feeds of 23~September had first been read as ``malformed XML'': the
feed library was parsing the HTML of the 404 page. Classifying the
outcome of the fetch before parsing it (Section~\ref{subsec:scraping-fetcher})
is what made the real cause visible.
